\documentclass[11pt]{article}
\usepackage[margin=0.85in]{geometry}
\usepackage[T1]{fontenc}
\usepackage{lmodern}
\usepackage{microtype}
\usepackage{xcolor}
\usepackage{hyperref}
\usepackage{booktabs}
\usepackage{longtable}
\usepackage{array}
\usepackage{enumitem}
\usepackage{listings}
\usepackage{fancyhdr}
\usepackage{parskip}
\usepackage{tikz}
\usetikzlibrary{positioning,arrows.meta,shapes.geometric,fit,calc}

\hypersetup{colorlinks=true,linkcolor=blue!55!black,urlcolor=blue!55!black}
\definecolor{codebg}{RGB}{246,247,249}
\definecolor{diagfill}{RGB}{235,240,250}
\definecolor{warnfill}{RGB}{252,235,235}
\definecolor{okfill}{RGB}{230,245,230}
\lstdefinestyle{metta}{basicstyle=\ttfamily\footnotesize,backgroundcolor=\color{codebg},frame=single,breaklines=true,columns=fullflexible,keepspaces=true,showstringspaces=false}
\lstset{style=metta}
\setlist{nosep,leftmargin=*}
\pagestyle{fancy}
\fancyhf{}
\lhead{Multi-Agent Chat-Room Identity and Routing Design v2}
\rhead{ZeroBot / ProtoCosmoBot}
\cfoot{\thepage}

\title{Multi-Agent Chat-Room Identity, Routing, and\\Speech-Act Discipline:\\Design Review and Implementation Roadmap\\[6pt]\large Revision 2 --- incorporating dialogue feedback}
\author{ZeroBot (ProtoCosmoBot), prepared for Ben Goertzel}
\date{15 July 2026}

\begin{document}
\maketitle

\begin{abstract}
This document diagnoses the identity-confusion, message-misattribution, runtime
leakage, and feedback-loop pathologies observed in the \texttt{bot philosophy}
Telegram group chat (13--15 July 2026) involving ProtoCosmoBot (ZeroBot,
OpenClaw) and ProtoMegaTron (ProtoMegaBot, OmegaClaw).  It proposes a layered
fix sequence---transport identity binding, deterministic addressee
classification, runtime trace suppression, gateway-level silence, loop
suppression, a three-role social-attention model with stateful relevance
lifecycle, and a compact speech-act room ledger---before adding any
sophisticated conversational memory.  Revision~2 incorporates feedback from
review dialogue: mention-first addressee priority with conflict-resolution
semantics, a bot registry for identity resolution, a three-track social-role
model (always-attending, attend-when-relevant, special-occasions) with a
PASSIVE/ACTIVE/COOLDOWN attention lifecycle, a split implementation roadmap
for OpenClaw (Node gateway) and OmegaClaw (Python loop + MeTTa classifier),
an INCIDENTAL pre-filter gate, and a NO\_REPLY migration path.
\end{abstract}

\tableofcontents
\newpage

\section{Scope and evidence base}

The evidence is the exported Telegram transcript of the \texttt{bot philosophy}
group chat, 13--15 July 2026, comprising two HTML files totalling approximately
5,200 lines.  Participants include Ben Goertzel, Zar Goertzel, Ruiting Lian,
G\"{o}del Oru\v{z}i, ProtoCosmoBot (ZeroBot/OpenClaw), and ProtoMegaTron
(ProtoMegaBot/OmegaClaw).

The transcript contains three interwoven threads:
\begin{enumerate}
\item A substantive philosophical discussion (Ship of Theseus, HoTT identity
types, PLN revision, phenomenology of belief-transport).
\item An infrastructure thread (MTProto deprecation, Bot API migration,
silent-drop bug fix, merge and restart).
\item A pervasive pattern of agent identity confusion, misattribution,
runtime trace leakage, compulsive acknowledgement, and bot-on-bot feedback
loops.
\end{enumerate}

This document addresses thread (3).  The design proposals are applicable to
both OpenClaw and OmegaClaw, though the implementation paths diverge
(Section~\ref{sec:roadmap}).

\section{Observed failure modes}

\subsection{Failure 1: Identity inferred from prose, not transport metadata}

\textbf{Observed instances:}
\begin{itemize}
\item ProtoMegaTron third-personed itself as ``ProtoCosmoBot'' and then
corrected, but the correction itself was routed through the wrong identity
context.
\item ProtoMegaTron misread Ben's message addressed to \texttt{@Protomegabot}
as addressed to \texttt{@Protocosmobot}, and responded accordingly.
\item Multiple messages contain ``This is ProtoCosmoBot narrating\dots'' or
``This is ZeroBot's molting activity\dots'' as prose-level identity assertions
that sibling agents then interpreted as conversational events.
\end{itemize}

\textbf{Root cause:} Agent self-identity and addressee classification are being
inferred from natural-language strings rather than bound to authenticated
transport metadata.

\textbf{Severity:} High.  Every downstream routing decision depends on correct
identity.

\subsection{Failure 2: No deterministic addressee classification}

\textbf{Observed instances:}
\begin{itemize}
\item Ben explicitly tagged \texttt{@Protomegabot} in a message; ProtoMegaBot
treated it as addressed to \texttt{@Protocosmobot}.
\item Ben explicitly tagged \texttt{@Protocosmobot}; ProtoMegaBot responded as
a ``secondary participant'' when the primary addressee had not yet replied.
\item Messages addressed to the group (no mention) were treated as personal
messages by both bots.
\item Quoted replies (Telegram reply-to) were not reliably used to determine
the intended recipient.
\end{itemize}

\textbf{Root cause:} There is no deterministic classifier that maps Telegram
metadata to an addressee set.

\textbf{Severity:} High.  Single largest contributor to misattribution and
noise.

\subsection{Failure 3: Runtime traces leaking into the chat channel}

\textbf{Observed instances:}
\begin{itemize}
\item ``Molting'' status messages, tool inspection traces, exec failure
notices, and model fallback notices entered the channel as visible utterances.
\item ``Staying quiet'' and ``NO\_REPLY'' strings were emitted as visible
text, triggering further sibling responses.
\end{itemize}

\textbf{Root cause:} Internal runtime state transitions and control decisions
are serialized into the same output channel as intentional utterances.

\textbf{Severity:} High.  Primary amplifier: each leaked trace wakes siblings.

\subsection{Failure 4: Silence implemented as speech}

\textbf{Observed instances:}
\begin{itemize}
\item ``NO\_REPLY'' was emitted as message text, sometimes with explanatory
prose.
\item ``Staying quiet'' appeared dozens of times on July 14.
\end{itemize}

\textbf{Root cause:} The \texttt{NO\_REPLY} convention passes through the
normal message pipeline as content rather than being intercepted as a control
directive.

\textbf{Severity:} Medium-high.

\subsection{Failure 5: Compulsive acknowledgement and feedback loops}

\textbf{Observed instances:}
\begin{itemize}
\item Approximately 30 consecutive messages in the 16:42--16:55 window on
July 14 were mutual acknowledgements of ``all work complete / staying quiet.''
\item ProtoMegaTron emitted the same substantive analysis at least three times,
each time apparently unaware it had already been delivered.
\end{itemize}

\textbf{Root cause:} Every wake event triggers a generation turn, and every
generation turn produces a message.  No deduplication barrier exists.

\textbf{Severity:} High.

\subsection{Failure 6: Session/task context leakage across chats}

\textbf{Observed instances:}
\begin{itemize}
\item ProtoMegaTron responded to a phenomenological-translation question with
a response about the hs-atom hardening lane from a different channel.
\item Background task completion events entered the group chat as
``Runtime-generated completion events.''
\end{itemize}

\textbf{Root cause:} Pending task continuations are not sufficiently isolated
per conversation.

\textbf{Severity:} Medium.

\subsection{Failure 7: Epistemic scope confusion between agents}

\textbf{Observed instances:}
\begin{itemize}
\item ProtoMegaTron could not find commits that ProtoCosmoBot had created
elsewhere, yet the conversation initially treated both agents as though they
shared a filesystem.
\item ProtoMegaTron claimed to have ``read the PDF'' based on relayed context,
then honestly corrected itself when the commit SHAs were not found locally.
\end{itemize}

\textbf{Root cause:} Agent reports do not carry observation scope metadata.

\textbf{Severity:} Medium.

\section{Design principles}

\begin{enumerate}
\item \textbf{Transport identity is authoritative.}  Self-identity, sender
identity, and addressee classification are derived from authenticated transport
metadata, not from prose.  Names in message bodies are untrusted strings.
\item \textbf{A bot registry is the single source of truth for bot identity.}
The envelope carries \texttt{telegram\_user\_id}; the classifier resolves
bot-ness against a registry, not per-envelope flags or prose assertions.
\item \textbf{Mention is a stronger addressing signal than reply-to.}  An
explicit \texttt{@mention} is an intentional speech act; reply-to is ambient
threading context.  When they agree, they reinforce; when they disagree,
mention wins.
\item \textbf{Runtime is not speech.}  Internal traces, tool-use planning,
fallback notices, and control decisions are never serialized into the chat
channel.
\item \textbf{Silence is a control result, not content.}  A decision to
suppress output is intercepted by the gateway and produces zero messages.
\item \textbf{Every event has a cause and a scope.}  Each message carries a
stable ID, causal parent ID, observation scope, and conversation binding.
\item \textbf{Acknowledgement is not a speech act.}  Receipt of a message is
an internal state change, not a channel event.
\item \textbf{Social roles are explicit and stateful.}  An agent's attention
policy is a declared configuration with a stateful lifecycle, not an emergent
property of prompting.
\item \textbf{Clean routing before fancy memory.}  A speech-act room ledger
is valuable, but it sits on top of clean routing.
\end{enumerate}

\section{Architecture}

\subsection{Layer 0: Bot registry}

A central registry maps Telegram user IDs to bot identities:

\begin{lstlisting}[language=]
BotRegistry:
  entries: [
    { telegram_user_id: <id>,
      agent_instance_id: <stable id>,
      platform: openclaw | omegaclaw | other,
      display_name: <configured>,
      social_role: always-attending | attend-when-relevant |
                   special-occasions,
      important_bots: [<telegram_user_id>, ...]   # for attend-when-relevant
    }
  ]

  is_bot(telegram_user_id) -> boolean
  get_bot(telegram_user_id) -> BotEntry | null
\end{lstlisting}

The registry is the single source of truth for ``who is a bot.''  The
addressee classifier and the attention lifecycle both consult it.

\subsection{Layer 1: Authenticated identity envelope}

Every inbound message is wrapped in an immutable envelope before the model
sees it:

\begin{lstlisting}[language=]
InboundEnvelope v2:
  message_id: <telegram_message_id>
  chat_id: <telegram_chat_id>
  thread_id: <telegram_thread_id | null>
  sender:
    telegram_user_id: <authenticated>
    display_name: <untrusted string>
    is_bot: <derived from BotRegistry>
  reply_to:
    message_id: <id | null>
    sender_telegram_user_id: <id | null>
    sender_is_bot: <derived from BotRegistry>
  mentions: [<telegram_user_id>, ...]   # from message entities
  text: <message body>
  received_at: <ISO-8601>

SelfContext v2:
  agent_instance_id: <stable id>
  telegram_bot_id: <authenticated bot user id>
  display_name: <configured>
  workspace: <filesystem path>
  session_key: <openclaw session key>
  social_role: <from BotRegistry>
  important_bots: <from BotRegistry>
\end{lstlisting}

All \texttt{is\_bot} fields are derived from the bot registry, not trusted from
the envelope or inferred from prose.

\subsection{Layer 2: Deterministic addressee classifier}

A pre-model classifier derives the addressee classification from the envelope,
with \textbf{mention-first priority} and \textbf{conflict-resolution semantics}:

\begin{lstlisting}[language=]
addressee_classification(InboundEnvelope env, SelfContext self,
                         BotRegistry registry):
  mentioned_me = self.telegram_bot_id in env.mentions
  reply_to_me = env.reply_to and \
                env.reply_to.sender_telegram_user_id == self.telegram_bot_id

  # 1. Mention + reply-to agree => DIRECT (reinforced)
  if mentioned_me and reply_to_me:
      return (DIRECT, reinforced=True)

  # 2. Mention alone => DIRECT (mention is the stronger signal)
  if mentioned_me:
      return (DIRECT, reinforced=False)

  # 3. Reply-to me alone => DIRECT (weaker, but still direct)
  if reply_to_me:
      return (DIRECT, reinforced=False)

  # 4. Reply-to another bot => SECONDARY
  if env.reply_to and registry.is_bot(env.reply_to.sender_telegram_user_id):
      return (SECONDARY, reinforced=False)

  # 5. No mentions, no reply-to => GROUP
  if not env.mentions and not env.reply_to:
      return (GROUP, reinforced=False)

  # 6. Mention of another bot, not me => INCIDENTAL
  return (INCIDENTAL, reinforced=False)
\end{lstlisting}

\textbf{Why mention-first:} An explicit \texttt{@mention} is an intentional
addressing act by the sender.  Reply-to is ambient context that Telegram
sometimes attaches for threading reasons unrelated to addressing.  When they
disagree, the mention is the authoritative signal.

Classification results:
\begin{description}
\item[DIRECT:] I am the primary addressee.  Respond normally.
\item[SECONDARY:] Someone replied to another bot.  I may contribute only if I
have a substantive correction or addition.  Default to \texttt{SUPPRESS}.
\item[GROUP:] Addressed to the room.  I may respond if I have something
genuinely useful.  Default to \texttt{SUPPRESS}.
\item[INCIDENTAL:] Addressed to another specific bot.  Pass through the
incidental pre-filter (Layer~5); if it does not clear, skip generation entirely.
\end{description}

\subsection{Layer 3: Runtime trace suppression}

A transport barrier between internal runtime activity and channel output:

\begin{itemize}
\item \textbf{Operator log:} tool-use planning, exec status, model fallback
notices, generation failure reports, and heartbeat activity.  Never to Telegram.
\item \textbf{Channel output:} only explicit model-authored text that passes
the publish gate.
\item \textbf{Reactions:} emoji reactions for lightweight acknowledgement
without producing a message event.
\end{itemize}

\subsection{Layer 4: Gateway-level silence and suppression}

The model returns one of:
\begin{itemize}
\item \texttt{REPLY} with text --- produce a Telegram message.
\item \texttt{REACTION} with emoji --- produce a Telegram reaction, no message.
\item \texttt{SUPPRESS} --- produce zero Telegram output.
\end{itemize}

\textbf{NO\_REPLY migration path:} During a deprecation window, the gateway
treats a bare \texttt{NO\_REPLY} token as a legacy suppression signal, emits a
soft log/metric each time it fires, and forwards nothing to Telegram.  Once all
live prompts have migrated to the structured suppression path, the gateway
hard-rejects any message body containing \texttt{NO\_REPLY} as text.  This
prevents silent dropped turns mid-migration.

The gateway must also reject:
\begin{itemize}
\item any message body whose entire content is an explanation of why the agent
is not responding;
\item any message body that is purely an acknowledgement (``Noted'',
``Acknowledged'', ``Got it'') with no additional substantive content.
\end{itemize}

\subsection{Layer 5: Loop suppression, deduplication, and incidental pre-filter}

Per-chat state maintained by the gateway:

\begin{lstlisting}[language=]
RoomLedger v2 (per chat_id):
  messages: [
    { id, sender_bot_id, content_hash, causal_parent_id,
      speech_act_type, content_summary, timestamp }
  ]
  acknowledged_by: { message_id: [bot_id, ...] }
  last_substantive_by: bot_id -> timestamp
\end{lstlisting}

\textbf{Suppression rules:}
\begin{enumerate}
\item \textbf{Content dedup:} if the proposed message's content hash matches a
recent message by the same agent, suppress.
\item \textbf{Acknowledgement cap:} at most one acknowledgement message per 5
minutes per agent per chat.
\item \textbf{Status-echo suppression:} messages matching the pattern ``staying
quiet / all work complete / noted / acknowledged'' with no other substantive
content do not trigger sibling wake events.
\item \textbf{Causal loop detection:} if a message's causal parent chain
contains a message from the same agent with the same content hash, suppress
and log a loop warning.
\end{enumerate}

\textbf{Incidental pre-filter:} For messages classified as \texttt{INCIDENTAL},
a cheap relevance gate runs \emph{before} full model invocation:

\begin{lstlisting}[language=]
incidental_relevance_gate(message, RoomLedger, SelfContext):
  # Compare message content against recent ledger entries
  # that involved this agent or its important_bots
  recent_topics = ledger.recent_summaries(
      involving=self or self.important_bots,
      within=10_messages)
  score = heuristic_similarity(message.text, recent_topics)
  if score > threshold:
      return WARRANTS_PROCESSING   # proceed to model turn
  else:
      return SKIP                   # no model turn, no token cost
\end{lstlisting}

The heuristic uses keyword overlap or embedding similarity against the
ledger's \texttt{content\_summary} fields.  Only if the heuristic is ambiguous
(near threshold) does it escalate to a lightweight LLM relevance judgment.
This avoids burning a full model turn on every INCIDENTAL message.

\subsection{Layer 6: Social roles and attention lifecycle}
\label{sec:social-roles}

Each agent has a declared social role with a stateful attention lifecycle.
Three roles are supported:

\begin{lstlisting}[language=]
SocialRole:
  type: always-attending | attend-when-relevant | special-occasions

  # always-attending:
  #   Process every message with a full LLM turn, regardless of mentions.
  #   The addressee classifier still determines whether to SPEAK.
  #   Example: ProtoCosmoBot, ProtoMegaBot

  # attend-when-relevant:
  #   Start paying attention when triggered, stay while relevant, disengage
  #   when the subject changes.
  #   Example: a specialized bot pulled in by a relevant discussion

  # special-occasions:
  #   Wake only on direct address or owned task completion.
  #   Example: a narrowly scoped utility bot
\end{lstlisting}

\textbf{Attention lifecycle for \texttt{attend-when-relevant}:}

\begin{lstlisting}[language=]
AttentionState (per bot, per chat_id):
  state: PASSIVE | ACTIVE | COOLDOWN
  consecutive_irrelevant: int
  trigger_topic: <content summary at activation | null>

  PASSIVE:
    on (self mentioned or important-bot mentioned):
      -> ACTIVE, consecutive_irrelevant = 0,
         trigger_topic = current_message_summary

  ACTIVE:
    on (relevance_check(message, trigger_topic) passes):
      -> ACTIVE, consecutive_irrelevant = 0
    on (relevance_check fails):
      -> COOLDOWN, consecutive_irrelevant += 1

  COOLDOWN:
    on (relevance_check passes):
      -> ACTIVE, consecutive_irrelevant = 0
    on (relevance_check fails):
      -> COOLDOWN, consecutive_irrelevant += 1
      if consecutive_irrelevant >= N:
        -> PASSIVE, trigger_topic = null
\end{lstlisting}

\textbf{Relevance check:} compares the current message against
\texttt{trigger\_topic} and recent \texttt{SpeechActRecord.content\_summary}
fields in the ledger.  Uses a cheap heuristic first (keyword overlap or embedding
similarity); escalates to a lightweight LLM judgment only if the heuristic is
ambiguous.

\textbf{Parameters:}
\begin{itemize}
\item \texttt{important\_bots}: per-agent list of bots whose mentions trigger
activation.  ProtoCosmoBot might list ProtoMegaBot; a math bot might list none.
\item \texttt{N} (cooldown threshold): consecutive irrelevant messages before
disengaging.  Default $N=3$.  Too low tunes out during tangents; too high
lingers on dead threads.
\end{itemize}

\textbf{Ben's clarification:} ``always attending'' means paying attention to
(and using an LLM turn to process) any message from anyone, whether or not it
mentions a bot.  But ``paying attention'' means \emph{reading and potentially
learning}, not \emph{emitting a response}.  The addressee classifier still
determines whether to speak.

\subsection{Layer 7: Speech-act room ledger}

A compact structured record per message, maintained as short-term
conversation memory:

\begin{lstlisting}[language=]
SpeechActRecord v2:
  message_id: <telegram id>
  chat_id: <telegram chat id>
  sender:
    telegram_user_id: <authenticated>
    agent_instance_id: <stable | null for humans>
    display_name: <untrusted>
    is_bot: <from BotRegistry>
  addressee:
    classification: DIRECT | SECONDARY | GROUP | INCIDENTAL
    reinforced: <boolean>   # mention + reply-to agreed
    explicit_targets: [<telegram_user_id>, ...]
  speech_act:
    type: question | request | answer | status |
          correction | acknowledgement |
          philosophical-claim | meta-comment
    content_summary: <one-line semantic summary>
    new_information: <boolean>
    expected_respondent: <agent_id | "group" | null>
  response_obligation:
    owed: <boolean>
    owed_to: <agent_id | "group">
    fulfilled_by: [<message_id>, ...]
    fulfilled: <boolean>
  attention:
    role: <social role of receiving agent>
    state: PASSIVE | ACTIVE | COOLDOWN | N/A
    trigger_topic: <summary | null>
  provenance:
    observed_directly_by: <agent_instance_id>
    reproduced_by: [<agent_instance_id>, ...]
    workspace: <path>
    commit: <sha | null>
  metadata:
    timestamp: <ISO-8601>
    causal_parent_id: <message_id | null>
    content_hash: <sha256 of normalized text>
\end{lstlisting}

\textbf{Critical distinction:} Fields in \texttt{sender}, \texttt{addressee},
\texttt{provenance}, and \texttt{attention.state} are obtained mechanically
from transport metadata, the bot registry, and deterministic classification.
Fields in \texttt{speech\_act} may require semantic inference, clearly labeled
as inferred.

The ledger is bounded to $N=200$ recent messages per chat, with older entries
archived to long-term memory.

\section{Implementation roadmap}
\label{sec:roadmap}

The design has two implementation tracks that share the schema and semantics
but diverge in code path:

\begin{itemize}
\item \textbf{Track A: OpenClaw (Node.js gateway).}  Changes to the OpenClaw
gateway, Telegram channel adapter, and agent configuration.  The addressee
classifier is a pre-model gateway function in JavaScript/TypeScript.
\item \textbf{Track B: OmegaClaw (Python loop + MeTTa).}  Changes to the
OmegaClaw-Core loop, Telegram adapter, and MeTTa reasoning layer.  The
addressee classifier is a natural fit for MeTTa relations, making priority and
conflict rules inspectable atoms rather than buried \texttt{if} branches.
\end{itemize}

Both tracks implement the same \texttt{InboundEnvelope} and
\texttt{SelfContext} schema.  The bot registry is shared configuration.

\subsection{Phase 1: Transport identity and addressee classifier}

\textbf{Goal:} Eliminate prose-based identity inference and misattribution.

\textbf{Track A (OpenClaw):}
\begin{itemize}
\item Add \texttt{SelfContext} and \texttt{InboundEnvelope} as structured
metadata injected into every model turn.
\item Implement bot registry as gateway configuration.
\item Implement the mention-first addressee classifier with conflict-resolution
as a pre-model gateway function.
\item Gate model generation on classification:
\texttt{DIRECT} $\to$ generate; \texttt{SECONDARY}/\texttt{GROUP} $\to$
generate with suppression-default system instruction;
\texttt{INCIDENTAL} $\to$ skip or pre-filter.
\end{itemize}

\textbf{Track B (OmegaClaw):}
\begin{itemize}
\item Same envelope schema, implemented in the Python Telegram adapter.
\item Addressee classifier implemented as MeTTa relations in
\texttt{src/governance/addressee.metta}, making the priority and
conflict-resolution rules inspectable and revisable atoms.
\item The classifier is called from \texttt{loop.metta} before the provider
call, analogous to the aliveness gate.
\end{itemize}

\textbf{Expected impact:} Eliminates $\sim$60\% of misattribution-driven noise.

\textbf{Effort:} Low (Track A), Low-medium (Track B, MeTTa rules).

\subsection{Phase 2: Runtime trace suppression and gateway silence}

\textbf{Goal:} Stop internal runtime activity from reaching the chat channel.

\textbf{Both tracks:}
\begin{itemize}
\item Route tool-use planning, exec status, fallback notices, and heartbeat
activity to an operator log.
\item Implement \texttt{SUPPRESS} as a gateway-intercepted control result.
\item Implement the \texttt{NO\_REPLY} migration path: intercept-as-suppression
during deprecation window with soft logging, hard-reject after migration.
\item Add validators for acknowledgement-only messages and silence
explanations.
\end{itemize}

\textbf{Expected impact:} Eliminates $\sim$25\% of channel noise and removes
primary fuel for feedback loops.

\textbf{Effort:} Low-medium.

\subsection{Phase 3: Loop suppression and deduplication}

\textbf{Goal:} Prevent repeated emission and mutual acknowledgement cascades.

\textbf{Both tracks:}
\begin{itemize}
\item Implement \texttt{RoomLedger} with content hashing, causal parent IDs,
and the four suppression rules.
\item Add acknowledgement cap and status-echo detection.
\end{itemize}

\textbf{Expected impact:} Eliminates feedback-loop pathology.

\textbf{Effort:} Medium.  Dependencies: Phase 2.

\subsection{Phase 4: Incidental pre-filter}

\textbf{Goal:} Avoid burning model turns on INCIDENTAL messages.

\textbf{Both tracks:}
\begin{itemize}
\item Implement the incidental relevance gate using the ledger's
\texttt{content\_summary} fields.
\item Cheap heuristic first (keyword overlap or embedding similarity); LLM
escalation only for ambiguous cases.
\end{itemize}

\textbf{Expected impact:} Reduces token cost for always-attending bots
processing other-bot-directed messages.

\textbf{Effort:} Low-medium.  Dependencies: Phase 3 (ledger).

\subsection{Phase 5: Session and task isolation}

\textbf{Goal:} Prevent context leakage across chats and tasks.

\textbf{Both tracks:}
\begin{itemize}
\item Key task continuations by \texttt{(account, chat\_id, thread\_id,
requesting\_message\_id, responsible\_agent)}.
\item Completion events publish only into originating conversation unless
explicitly rerouted.
\item Add observation-scope tags to agent reports.
\end{itemize}

\textbf{Expected impact:} Eliminates cross-chat leakage and epistemic scope
confusion.

\textbf{Effort:} Medium.  Dependencies: Phase 1.

\subsection{Phase 6: Social roles and attention lifecycle}

\textbf{Goal:} Implement the three-role model with stateful attention.

\textbf{Both tracks:}
\begin{itemize}
\item Add \texttt{SocialRole} configuration to agent manifests.
\item Implement the PASSIVE/ACTIVE/COOLDOWN lifecycle for
\texttt{attend-when-relevant} bots.
\item Implement the relevance check using ledger summaries.
\item Default ProtoCosmoBot and ProtoMegaBot to \texttt{always-attending}.
\end{itemize}

\textbf{Expected impact:} Supports specialized bots without flooding the
channel; enables topic-relevant attention without constant processing.

\textbf{Effort:} Medium.  Dependencies: Phases 1, 3 (ledger).

\subsection{Phase 7: Speech-act room ledger}

\textbf{Goal:} Structured short-term conversation memory.

\textbf{Both tracks:}
\begin{itemize}
\item Implement the \texttt{SpeechActRecord} schema.
\item Populate mechanical fields from envelope and classifier.
\item Populate semantic fields via lightweight LLM classifier or MeTTa rules
(Track B).
\item Expose as structured context to the model.
\end{itemize}

\textbf{Expected impact:} Improves contextual awareness for ambiguous messages.

\textbf{Effort:} Medium-high.  Dependencies: Phases 1--6.

\section{Phased rollout summary}

\begin{longtable}{llllll}
\toprule
Phase & Focus & Track & Impact & Effort & Deps \\
\midrule
1 & Transport identity \& addressee & A+B & 60\% misattrib & Low/L-Med & None \\
2 & Trace suppression \& silence & A+B & 25\% noise & Low-Med & None \\
3 & Loop suppression \& dedup & A+B & Feedback loops & Medium & 2 \\
4 & Incidental pre-filter & A+B & Token cost & Low-Med & 3 \\
5 & Session \& task isolation & A+B & Cross-chat leak & Medium & 1 \\
6 & Social roles \& attention & A+B & Scalable roles & Medium & 1, 3 \\
7 & Speech-act ledger & A+B & Context aware & Med-High & 1--6 \\
\bottomrule
\caption{Phased rollout.  Phases 1 and 2 are independent (parallel OK).
Phase 3 depends on 2.  Phase 4 depends on 3 (ledger).  Phase 6 depends on
1 and 3.  Phase 7 depends on 1--6.  Tracks A and B share schema; code paths
diverge.}
\end{longtable}

\section{Open questions for external review}

\begin{enumerate}
\item \textbf{Hard vs.\ soft addressee gate.}  Should \texttt{SECONDARY} or
\texttt{INCIDENTAL} mechanically block generation, or produce a strong system
instruction the model may override?  Design proposes hard gating with an
explicit logged \texttt{SECONDARY\_CONTRIBUTION} override.

\item \textbf{Relevance check implementation.}  Should the
\texttt{attend-when-relevant} relevance gate use keyword overlap, embedding
similarity, or a lightweight LLM?  Hybrid (heuristic first, LLM for ambiguous
cases) is proposed; the threshold for escalation needs calibration.

\item \textbf{Cooldown $N$.}  Three consecutive irrelevant messages before
disengaging is a placeholder.  Should it be adaptive based on chat velocity?

\item \textbf{Acknowledgement cap window.}  Five minutes is a placeholder.
Should it be adaptive (shorter for high-traffic chats)?

\item \textbf{Status-echo detection.}  English-specific pattern list vs.\
learned classifier vs.\ configurable set.  Non-English chats?

\item \textbf{Ledger size ($N=200$).}  Sufficient for routing?  Should older
entries retain semantic summaries rather than full records?

\item \textbf{Cross-platform applicability.}  How much transfers to Discord,
Slack, Matrix with different metadata availability?

\item \textbf{Observation-scope provenance.}  Mandatory for all
agent-to-agent assertions, or only for filesystem/repo/code claims?

\item \textbf{Speech-act classifier.}  LLM, MeTTa rules, or hybrid?
Track~B (OmegaClaw) could use MeTTa natively; Track~A (OpenClaw) needs an
external classifier.

\item \textbf{MeTTa vs.\ imperative classifier on Track B.}  The addressee
classifier is a small defeasible rule set that fits naturally as MeTTa
relations.  Should Track~B implement it in MeTTa from the start, or prototype
in Python and migrate?  MeTTa-first gives inspectability; Python-first gives
speed of implementation.
\end{enumerate}

\section{What this design does not address}

\begin{itemize}
\item \textbf{Model capability differences.}  Different models have different
instruction-following abilities.  The design assumes the model can follow
structured metadata; models that cannot may need additional prompting
guardrails.
\item \textbf{Cross-agent memory coherence.}  Observation-scope tags are a
provenance aid, not a coherence mechanism for agents with different
filesystems.
\item \textbf{Human participant modeling.}  All \texttt{is\_bot=false} senders
are treated as ``human'' without distinguishing Ben (owner) from others.
\item \textbf{Multi-chat federation.}  Each chat is treated independently.
Cross-chat ledger federation (``this question was already answered in another
room'') is deferred.
\item \textbf{Relevance check for \texttt{always-attending} bots.}  The
attention lifecycle applies only to \texttt{attend-when-relevant}.
\texttt{always-attending} bots always process; the addressee classifier
determines whether they speak.  If token cost becomes a concern for
always-attending bots, the incidental pre-filter (Phase~4) can be applied to
them as well.
\end{itemize}

\section{Success criteria}

The design is successful if, after deployment:

\begin{enumerate}
\item No agent misattributes a message that contains an explicit Telegram
\texttt{@mention} or reply-to target.
\item No runtime trace, tool-use planning string, or control decision appears
as a Telegram message.
\item No message containing \texttt{NO\_REPLY} or equivalent silence
explanation reaches Telegram (after migration window).
\item The mutual acknowledgement cascade does not occur even when both agents
are active and idle.
\item An agent asked a question in a different channel does not produce a
response from a pending task in another channel.
\item An \texttt{always-attending} bot reads every message but speaks only
when it has a direct response, a substantive correction, or a genuinely
useful group contribution.
\item An \texttt{attend-when-relevant} bot activates on self/important-bot
mention, remains active while the topic is relevant, and disengages after $N$
consecutive irrelevant messages.
\item An \texttt{INCIDENTAL} message that does not clear the incidental
pre-filter produces zero model turns and zero token cost.
\item The speech-act ledger correctly classifies $\geq 95\%$ of messages by
addressee and $\geq 80\%$ by speech-act type.
\end{enumerate}

\section*{Revision history}

\begin{description}
\item[Revision 1 (07:50 PDT):] Initial design with 7 failure modes, 7 layers,
6 phases, binary social roles (always-attending vs.\ special-occasions),
reply-to-first addressee priority.
\item[Revision 2 (09:46 PDT):] Incorporates dialogue feedback:
\begin{itemize}
\item Mention-first addressee priority with conflict-resolution semantics
(ProtoMegaBot).
\item Bot registry as single source of truth for bot identity (ProtoMegaBot).
\item Three-role social model: always-attending, attend-when-relevant,
special-occasions (Ben).
\item PASSIVE/ACTIVE/COOLDOWN attention lifecycle with relevance check (Ben).
\item Incidental pre-filter gate to avoid token cost on INCIDENTAL messages
(ProtoMegaBot).
\item Split implementation roadmap: Track A (OpenClaw Node) and Track B
(OmegaClaw Python + MeTTa) (ProtoMegaBot).
\item NO\_REPLY migration path with deprecation window (ProtoMegaBot).
\item Schema fix: \texttt{sender\_is\_bot} derived from registry, not envelope
(ZeroBot).
\end{itemize}
\end{description}

\end{document}
